% !TEX root = ../main.tex

\chapter{相关技术}


\section{HTLC 哈希时间锁}


HTLC 全称为 Hashed TimeLock Contract，即哈希时间锁合约~\cite{poon2016bitcoin,10.1145/3212734.3212736}。它被广泛应用于链下支付通道和跨链原子交换中。其核心思想是将资产领取条件拆分为两个部分：哈希锁和时间锁。哈希锁要求领取方必须提交一个 secret，使得该 secret 经过哈希计算后等于合约中预先保存的 hashLock；时间锁则规定如果在指定时间内没有人提交正确 secret，资产可以被原始锁定方取回。


HTLC 的基本流程可以描述为：Alice 首先生成随机 secret，并计算 hashLock = H(secret)。在发起跨链交换时，Alice 并不公开 secret，而是将 hashLock 写入源链锁定合约。目标链也使用同一个 hashLock 创建对应锁定。这样，两条链上的资产虽然位于不同账本中，但它们的领取条件都被同一个 secret 绑定。当某一方为了领取目标链资产而公开 secret 时，另一方也可以使用该 secret 去完成源链或目标链上的对应操作。若 secret 始终没有公开，则双方都无法领取对方资产，最终只能等待 timeLock 到期后退款。


HTLC 的优势在于结构清晰、其原子性保持能力不依赖于对区块链账本状态的掌握，并且适合部署在具有智能合约能力的链上。在 Solidity 合约中，hashLock 可以作为 bytes32 类型保存，timeLock 可以用区块高度或时间戳表示。完成交易时，合约只需验证 keccak256(secret) 是否等于 hashLock；退款时，合约只需判断当前区块高度是否超过 timeLock。通过这种方式，链上合约能够自动执行“正确 secret 领取”和“超时退款”两类规则。


不过，HTLC 也存在适用边界~\cite{tsabary2021mad}。它要求参与链能够在链上表达哈希锁判断和超时退款两类条件逻辑，并要求参与方在超时前后及时发起领取或退款交易。这一前提在具备脚本或合约能力的链上容易满足——例如 BTC 通过 OP\_SHA256、OP\_HASH160 和 OP\_CHECKLOCKTIMEVERIFY 等操作码即可实现链上 HTLC，Lightning Network 正是建立在这一能力之上~\cite{poon2016bitcoin}。但对于 Monero、ZCash 屏蔽交易等无脚本链，情况则完全不同：这类链以环签名、零知识证明等纯密码学方式构造交易，不提供通用脚本系统，无法在链上表达``提交正确 secret 才能领取"这样的条件分支。因此，HTLC 的条件逻辑根本无法迁移到无脚本链上。要让这类链参与原子交换，必须将条件锁定从脚本层转移到签名协议层，本文选用在下一节所介绍的 Schnorr adaptor signature~\cite{9833731} 算法解决以上问题。


\section{Schnorr 签名与 Adaptor Signature}


\subsection{Schnorr 签名基础}


Schnorr 签名是一种基于离散对数问题的数字签名方案~\cite{schnorr1991efficient}，其核心特点是结构简洁且具有线性性质。设椭圆曲线群的生成元为 \(G\)，签名者的私钥为 \(x\)，对应公钥为 \(P = xG\)。对消息 \(m\) 的签名过程如下：签名者选取随机数 \(k\)，计算承诺点 \(R = kG\)，然后计算挑战值 \(e = H(R, P, m)\)，最后计算签名标量 \(s = k + e \cdot x\)。最终签名为 \((R, s)\)。


验证者收到签名 \((R, s)\) 后，计算 \(e = H(R, P, m)\)，并检验等式 \(sG = R + e \cdot P\) 是否成立。由于 \(sG = (k + e \cdot x)G = kG + e \cdot xG = R + e \cdot P\)，合法签名必然通过验证。


Schnorr 签名的线性性质是 adaptor signature 的数学基础。由于签名标量 s 是随机数 k 与私钥 x 的线性组合，可以在 s 上叠加或拆分标量值而不破坏签名的代数结构。这一性质使得”将秘密值嵌入签名补全过程”成为可能。


\subsection{Adaptor Signature 构造}


Adaptor signature（适配签名）~\cite{9833731} 在 Schnorr 签名的基础上引入一个额外的秘密点 \(T = tG\)，其中 \(t\) 是待传递的秘密标量。适配签名协议包含四个阶段：预签名（PreSign）、预验证（PreVerify）、补全（Adapt）和提取（Extract）。


PreSign 阶段：签名者选取随机数 \(k\)，计算调整后的承诺点 \(R' = kG + T\)，然后计算挑战值 \(e = H(R', P, m)\) 和预签名标量 \(s' = k + e \cdot x\)。输出预签名 \((R', s')\)。注意此时 \(s'\) 并不是相对于 \(R'\) 的有效 Schnorr 签名，因为 \(s'G = kG + e \cdot P = R' - T + e \cdot P\)，验证等式 \(s'G = R' + e \cdot P\) 不成立。


PreVerify 阶段：验证者已知公开的秘密点 \(T\)，可以检验 \(s'G = R' - T + e \cdot P\) 是否成立。该等式验证了预签名的正确性——即确认签名者确实使用了与 \(T\) 对应的承诺偏移，且签名者知道私钥 \(x\)。


Adapt 阶段：持有秘密 \(t\) 的一方计算完整签名标量 \(s = s' + t\)。此时 \(sG = s'G + tG = (R' - T + e \cdot P) + T = R' + e \cdot P\)，验证等式成立，\((R', s)\) 构成对消息 \(m\) 的有效 Schnorr 签名。该签名可以被广播上链，完成交易。


Extract 阶段：另一方观察到链上出现的有效签名 \(s\) 和此前收到的预签名 \(s'\)，通过计算 \(t = s - s'\) 即可提取秘密值。由于 \(s\) 和 \(s'\) 都是公开可见的标量，秘密提取是确定性的。


\subsection{密码学性质}


Adaptor signature 具有以下关键性质，使其适用于原子交换场景：


完整性：如果预签名正确生成且秘密 t 正确补入，补全后的签名必然是有效的 Schnorr 签名，能够通过链上验证。


不可伪造性：不知道私钥 x 的攻击者无法生成能通过 PreVerify 的预签名，不知道秘密 t 的攻击者无法将预签名补全为有效签名。


秘密提取的必然性：一旦补全后的签名被广播上链，任何持有预签名的观察者都能通过 s - s' 提取秘密 t。交易完成与秘密公开之间存在密码学上的必然绑定——不可能在不暴露 t 的情况下使用补全后的签名。


这三条性质共同保证了：在跨链交换中，一方完成交易（广播有效签名）必然导致秘密被另一方获得，从而使另一方也具备完成其侧交易的条件。这与 HTLC 中”公开 secret 领取资产”的逻辑等价，但条件绑定发生在签名协议层而非链上脚本层，链上只呈现一笔普通 Schnorr 签名交易。
